\documentclass[UTF8,a4paper,11pt]{ctexart}

\usepackage[margin=2.2cm]{geometry}
\usepackage{xcolor}
\usepackage{hyperref}
\usepackage{enumitem}
\usepackage{longtable}
\usepackage{booktabs}
\usepackage{listings}
\usepackage{caption}

\definecolor{codebg}{RGB}{248,248,248}
\definecolor{codeframe}{RGB}{220,220,220}
\definecolor{keywordblue}{RGB}{0,78,150}
\definecolor{commentgreen}{RGB}{0,120,70}
\definecolor{stringred}{RGB}{160,40,40}

\lstdefinestyle{codestyle}{
  basicstyle=\ttfamily\small,
  backgroundcolor=\color{codebg},
  frame=single,
  rulecolor=\color{codeframe},
  breaklines=true,
  tabsize=2,
  showstringspaces=false,
  keywordstyle=\color{keywordblue}\bfseries,
  commentstyle=\color{commentgreen},
  stringstyle=\color{stringred}
}
\lstset{style=codestyle}

\title{双臂机器人上位机/下位机代码按优先级改造清单}
\author{Codex 静态代码审阅整理}
\date{2026年5月19日}

\begin{document}
\maketitle

\section*{说明}

本文基于当前工程静态阅读整理，目标是把现有单臂控制代码可靠扩展为双臂控制系统。优先级定义如下：

\begin{itemize}[leftmargin=2em]
  \item \textbf{P0：必须先修。} 存在越界、错误下发、急停失效或会直接导致双臂不可控的问题。
  \item \textbf{P1：双臂主链路必需。} 不修也许能单臂运行，但无法稳定支持左右臂独立控制。
  \item \textbf{P2：可维护性和调试效率。} 不一定立刻阻塞实机，但会显著增加调试成本。
  \item \textbf{P3：体验和工程化完善。} 适合在主链路稳定后补齐。
\end{itemize}

\section{总体结论}

当前项目已经具备上位机 ROS2 控制循环、状态机、KDL 运动学/动力学、USB 串口通信、下位机 USB CDC 收发、FDCAN 下发和夹爪 PWM 控制等基础能力。主要风险集中在：下位机接收缓冲区长度不一致、右臂 CAN 下发复用了左臂数据、上位机控制层大量接口仍按 6 轴单臂编写、命令入口和轨迹来源未统一、诊断与安全闭环不足。

\section{P0：必须先修的问题}

\subsection{P0-1 修正下位机 USB 接收数据缓存越界}

\textbf{现象：} \verb|crc16.c| 中定义 \verb|USB_Rcive_Data[64]|，但帧数据长度已经是 112 字节；\verb|main.c| 又声明为 117 字节并访问到 \verb|USB_Rcive_Data[104+i]|。这会造成内存越界，属于实机运行前必须修复的问题。

\textbf{修改位置：}
\begin{itemize}[leftmargin=2em]
  \item \verb|double_arm_h7/vofa+/crc16.c|
  \item \verb|double_arm_h7/Core/Src/main.c|
\end{itemize}

\textbf{建议代码：}
\begin{lstlisting}[language=C]
// double_arm_h7/vofa+/crc16.c
// 原代码:
// volatile uint8_t USB_Rcive_Data[64];

#include "crc16.h"

volatile uint8_t USB_Rcive_Data[FRAME_DATA_LENGTH] = {0};
\end{lstlisting}

\begin{lstlisting}[language=C]
// double_arm_h7/Core/Src/main.c
// 原代码:
// extern volatile uint8_t USB_Rcive_Data[117];

#include "crc16.h"

extern volatile uint8_t USB_Rcive_Data[FRAME_DATA_LENGTH];
\end{lstlisting}

\textbf{验证方式：} 编译下位机工程，确认无重复定义/数组大小冲突；上位机持续发送 117 字节帧，观察下位机不 HardFault，左右臂 14 个 8 字节槽位均能正确更新。

\subsection{P0-2 修正右臂 CAN 下发复用左臂控制数据}

\textbf{现象：} 下位机主循环中右臂 FDCAN2 当前仍下发 \verb|motor1_date| 到 \verb|motor6_date|，没有使用 \verb|motor8_date| 到 \verb|motor13_date|。这会导致右臂收到左臂指令，是双臂实机的高风险问题。

\textbf{修改位置：} \verb|double_arm_h7/Core/Src/main.c|

\textbf{建议代码：}
\begin{lstlisting}[language=C]
// 左臂: FDCAN1, 电机 1-6
fdcanx_send_data(&hfdcan1, 0x01, motor1_date, 8);
fdcanx_send_data(&hfdcan1, 0x02, motor2_date, 8);
fdcanx_send_data(&hfdcan1, 0x03, motor3_date, 8);
fdcanx_send_data(&hfdcan1, 0x04, motor4_date, 8);
fdcanx_send_data(&hfdcan1, 0x05, motor5_date, 8);
fdcanx_send_data(&hfdcan1, 0x06, motor6_date, 8);

// 右臂: FDCAN2, 电机 8-13
fdcanx_send_data(&hfdcan2, 0x01, motor8_date, 8);
fdcanx_send_data(&hfdcan2, 0x02, motor9_date, 8);
fdcanx_send_data(&hfdcan2, 0x03, motor10_date, 8);
fdcanx_send_data(&hfdcan2, 0x04, motor11_date, 8);
fdcanx_send_data(&hfdcan2, 0x05, motor12_date, 8);
fdcanx_send_data(&hfdcan2, 0x06, motor13_date, 8);
\end{lstlisting}

\textbf{验证方式：} 上位机分别只改变左臂/右臂某一关节指令，使用 CAN 分析仪或电机反馈确认只有对应总线对应电机发生变化。

\subsection{P0-3 让实时保护真正急停位置越界}

\textbf{现象：} 上位机 \verb|real_time_motor_protection| 中速度和力矩越界会调用 \verb|DisableMotors|，但位置越界分支里失能动作被注释。机械臂位置越界时只打印警告，不会真正关断电机。

\textbf{修改位置：} \verb|ros2_ws_dbarm/src/light_lift_arm_6dof/src/drivers/motor_control.cpp|

\textbf{建议代码：}
\begin{lstlisting}[language=C++]
// 建议抽出统一急停函数，避免多个分支逻辑不一致。
void MotorControl::emergencyStop(serialport::SerialPortWrapper &port,
                                 const std::string &reason,
                                 int motor_index)
{
    std::cerr << "[EMERGENCY STOP] motor " << motor_index
              << ": " << reason << std::endl;
    for (size_t retry = 0; retry < 10; ++retry) {
        DisableMotors(port);
        std::this_thread::sleep_for(std::chrono::milliseconds(10));
    }
}

// 位置越界分支中调用:
if (pos_dengerous_cout[i] >= pos_max_dengerous_cout) {
    emergencyStop(port, "position out of range", static_cast<int>(i));
    bad_motor[i] = 1;
    out_flag = true;
}
\end{lstlisting}

\textbf{同步修改头文件：}
\begin{lstlisting}[language=C++]
// motor_control.h
void emergencyStop(serialport::SerialPortWrapper &port,
                   const std::string &reason,
                   int motor_index);
\end{lstlisting}

\textbf{验证方式：} 人为设置很小的位置保护上限，运行后确认触发后电机被失能，而不是只打印日志。

\subsection{P0-4 修正状态切换后的局部初始化状态没有回写}

\textbf{现象：} \verb|ArmState::run| 中把 \verb|is_init_end| 复制到局部变量 \verb|initialized|，但在切换模式时没有统一复位初始化标志。结果是部分模式切换后可能沿用旧目标/旧增益/旧期望，导致突跳。

\textbf{修改位置：} \verb|ros2_ws_dbarm/src/light_lift_arm_6dof/src/controllers/state_machine.cpp|

\textbf{建议代码：}
\begin{lstlisting}[language=C++]
void ArmState::modeTransition(State newState)
{
    if (currentState == newState) {
        return;
    }

    currentState = newState;

    // 模式切换时强制重新初始化目标，避免沿用上一模式的期望值。
    is_init_end = false;
    is_init_joint = false;
    is_run_planning = false;
    is_go_home = false;

    if (newState == State::Zero) {
        is_init_end = true;
        is_init_joint = true;
    }
    if (newState == State::Gohome) {
        is_go_home = true;
    }
}
\end{lstlisting}

\textbf{验证方式：} 在 AutoServo、Gravity、Impdence、JointAutoServo 之间多次切换，确认第一次进入新模式时会重新初始化期望位置，不出现旧轨迹突跳。

\section{P1：双臂主链路必需改造}

\subsection{P1-1 上位机电机控制接口从 6 轴扩展为 12 轴 + 2 夹爪}

\textbf{现象：} 下位机协议已经预留 112 字节，适合 14 个 8 字节槽位；但上位机 \verb|ControlMotors| 仍主要按 6 个电机打包，右臂数据没有完整填入。

\textbf{修改位置：}
\begin{itemize}[leftmargin=2em]
  \item \verb|motor_control.h|
  \item \verb|motor_control.cpp|
\end{itemize}

\textbf{建议代码：}
\begin{lstlisting}[language=C++]
// motor_control.h
static constexpr size_t ARM_DOF = 6;
static constexpr size_t ARM_COUNT = 2;
static constexpr size_t MOTOR_NUM = ARM_DOF * ARM_COUNT;
static constexpr size_t SLOT_SIZE = 8;
static constexpr size_t LEFT_GRIPPER_SLOT = 6;   // bytes 48-55
static constexpr size_t RIGHT_GRIPPER_SLOT = 13; // bytes 104-111

struct ArmCommand {
    float pos[ARM_DOF] = {0};
    float vel[ARM_DOF] = {0};
    float kp[ARM_DOF] = {0};
    float kd[ARM_DOF] = {0};
    float tau[ARM_DOF] = {0};
    uint8_t gripper = 0;
};

void ControlDualArm(serialport::SerialPortWrapper &port,
                    const ArmCommand &left,
                    const ArmCommand &right);
\end{lstlisting}

\begin{lstlisting}[language=C++]
// motor_control.cpp
static size_t commandOffset(size_t motor_index)
{
    size_t offset = motor_index * SLOT_SIZE;
    if (motor_index >= ARM_DOF) {
        offset += SLOT_SIZE; // skip left gripper slot
    }
    return offset;
}

void MotorControl::ControlDualArm(serialport::SerialPortWrapper &port,
                                  const ArmCommand &left,
                                  const ArmCommand &right)
{
    std::vector<uint8_t> data(port.FRAME_DATA_LENGTH, 0);
    uint8_t motor_data[8] = {0};

    const ArmCommand arms[2] = {left, right};

    for (size_t arm = 0; arm < ARM_COUNT; ++arm) {
        for (size_t joint = 0; joint < ARM_DOF; ++joint) {
            const size_t global = arm * ARM_DOF + joint;
            if (joint < 3) {
                MitCtrl4340(arms[arm].pos[joint], arms[arm].vel[joint],
                            arms[arm].kp[joint], arms[arm].kd[joint],
                            arms[arm].tau[joint], motor_data);
            } else {
                MitCtrl(arms[arm].pos[joint], arms[arm].vel[joint],
                        arms[arm].kp[joint], arms[arm].kd[joint],
                        arms[arm].tau[joint], motor_data);
            }

            const size_t offset = commandOffset(global);
            std::copy(motor_data, motor_data + SLOT_SIZE, data.begin() + offset);
        }
    }

    data[LEFT_GRIPPER_SLOT * SLOT_SIZE] = left.gripper;
    data[RIGHT_GRIPPER_SLOT * SLOT_SIZE] = right.gripper;

    std::vector<uint8_t> frame;
    port.packFrame(data, frame);
    port.sendData(frame);
}
\end{lstlisting}

\textbf{验证方式：} 构造左右臂不同位置命令，打印 112 字节 payload，确认左臂在 0--47，左夹爪在 48，右臂在 56--103，右夹爪在 104。

\subsection{P1-2 抽象左右臂状态，避免所有控制器直接依赖单个 \texttt{LlArm6dof}}

\textbf{现象：} 当前 KDL 模型只取 \verb|base_link -> link6| 单链，状态机也只维护一套当前关节和期望关节。要做双臂，建议先把左右臂状态封装成两个实例。

\textbf{修改位置：}
\begin{itemize}[leftmargin=2em]
  \item 新增 \verb|dual_arm_system.h/.cpp|
  \item 改造 \verb|light_lift_arm_6dof_node.cpp|
\end{itemize}

\textbf{建议代码：}
\begin{lstlisting}[language=C++]
enum class ArmSide { Left, Right };

struct ArmRuntime {
    LlArm6dof model;
    ArmState state_machine;
    MotorControl::ArmCommand command;
    std::string ns;
};

class DualArmSystem {
public:
    ArmRuntime left{"left"};
    ArmRuntime right{"right"};

    ArmRuntime& arm(ArmSide side)
    {
        return side == ArmSide::Left ? left : right;
    }

    void updateFromMotorState(const MotorControl &motor)
    {
        left.model.updatArmState(motor.current_motor_pos,
                                 motor.current_motor_vel,
                                 motor.current_motor_tau);

        right.model.updatArmState(motor.current_motor_pos + 6,
                                  motor.current_motor_vel + 6,
                                  motor.current_motor_tau + 6);
    }
};
\end{lstlisting}

\textbf{验证方式：} 上位机收到 12 电机反馈后，左臂 joint state 使用索引 0--5，右臂 joint state 使用索引 6--11，互不覆盖。

\subsection{P1-3 统一 ROS2 话题命名空间}

\textbf{现象：} 当前话题只有一套，例如 \verb|topic_control_desir_joint_state|。双臂需要左右臂独立控制，同时保留一个双臂协调命令入口。

\textbf{建议话题：}
\begin{longtable}{lll}
\toprule
用途 & 左臂 & 右臂 \\
\midrule
当前关节 & \verb|/left_arm/joint_states| & \verb|/right_arm/joint_states| \\
期望关节 & \verb|/left_arm/desir_joint_states| & \verb|/right_arm/desir_joint_states| \\
末端目标 & \verb|/left_arm/target_pose| & \verb|/right_arm/target_pose| \\
命令入口 & \verb|/left_arm/cmd| & \verb|/right_arm/cmd| \\
\bottomrule
\end{longtable}

\textbf{建议代码：}
\begin{lstlisting}[language=C++]
auto makeJointPublisher =
    [this](const std::string &name) {
        return this->create_publisher<sensor_msgs::msg::JointState>(
            name + "/joint_states", 10);
    };

left_joint_pub_ = makeJointPublisher("/left_arm");
right_joint_pub_ = makeJointPublisher("/right_arm");

left_pose_sub_ = this->create_subscription<geometry_msgs::msg::Pose>(
    "/left_arm/target_pose", 10,
    [this](geometry_msgs::msg::Pose::SharedPtr msg) {
        handleTargetPose(ArmSide::Left, *msg);
    });

right_pose_sub_ = this->create_subscription<geometry_msgs::msg::Pose>(
    "/right_arm/target_pose", 10,
    [this](geometry_msgs::msg::Pose::SharedPtr msg) {
        handleTargetPose(ArmSide::Right, *msg);
    });
\end{lstlisting}

\textbf{验证方式：} 使用 \verb|ros2 topic list| 确认左右臂话题分离；分别发布左右臂目标，确认只影响对应臂。

\subsection{P1-4 统一命令解析，补齐全部状态入口}

\textbf{现象：} 当前 \verb|cmd_callback| 只处理 Gravity、Impdence、JointAutoServo。建议把字符串映射集中管理，避免模式名拼写错误和漏分支。

\textbf{修改位置：} \verb|light_lift_arm_6dof_node.cpp|

\textbf{建议代码：}
\begin{lstlisting}[language=C++]
std::optional<State> parseStateCommand(const std::string &cmd)
{
    static const std::unordered_map<std::string, State> states = {
        {"AutoServo", State::AutoServo},
        {"ManualServo", State::ManualServo},
        {"Impdence", State::Impdence},
        {"Admittance", State::Admittance},
        {"Gravity", State::Gravity},
        {"JointAutoServo", State::JointAutoServo},
        {"Zero", State::Zero},
        {"Planning", State::Planning},
        {"Gohome", State::Gohome},
    };

    auto it = states.find(cmd);
    if (it == states.end()) {
        return std::nullopt;
    }
    return it->second;
}

void cmd_callback(const std_msgs::msg::String::SharedPtr msg)
{
    llarm6dof.cmd = msg->data;

    if (auto state = parseStateCommand(msg->data)) {
        mode_ = *state;
        control_mode.modeTransition(mode_);
        return;
    }

    if (msg->data == "load" || msg->data == "clearload") {
        return; // 负载估计命令由阻抗控制分支处理
    }

    RCLCPP_WARN(this->get_logger(), "Unknown command: %s", msg->data.c_str());
}
\end{lstlisting}

\textbf{验证方式：} 对每个模式发布一次命令，确认日志中的状态和实际运行分支一致。

\subsection{P1-5 明确轨迹目标来源，避免内部轨迹和话题目标互相覆盖}

\textbf{现象：} 主循环无条件调用 \verb|update_end_effector_trajectory|，而注释要求话题控制时必须关闭它。建议用枚举明确目标来源，而不是靠手动注释。

\textbf{建议代码：}
\begin{lstlisting}[language=C++]
enum class TargetSource {
    InternalTrajectory,
    TopicPose,
    TopicJoint,
    Planner
};

TargetSource target_source_ = TargetSource::InternalTrajectory;

void setTargetSource(TargetSource source)
{
    target_source_ = source;
}

void updateTarget(float realtime)
{
    switch (target_source_) {
    case TargetSource::InternalTrajectory:
        llarm6dof.update_end_effector_trajectory(realtime);
        break;
    case TargetSource::TopicPose:
    case TargetSource::TopicJoint:
    case TargetSource::Planner:
        break; // 已由回调或规划器更新目标
    }
}
\end{lstlisting}

\textbf{验证方式：} 切换话题控制后，内部轨迹不再覆盖目标；切回内部轨迹后，目标由函数更新。

\section{P2：可维护性和调试效率}

\subsection{P2-1 发布末端状态和 TF，补齐观测闭环}

\textbf{现象：} 当前已经创建了末端位姿、速度、力话题和 TF broadcaster，但发布语句大多被注释。建议启用并按左右臂命名。

\textbf{建议代码：}
\begin{lstlisting}[language=C++]
void publishEndEffectorState(const LlArm6dof &arm,
                             rclcpp::Publisher<geometry_msgs::msg::Pose>::SharedPtr pose_pub,
                             rclcpp::Publisher<geometry_msgs::msg::Twist>::SharedPtr twist_pub,
                             rclcpp::Publisher<geometry_msgs::msg::Wrench>::SharedPtr wrench_pub)
{
    geometry_msgs::msg::Pose pose;
    pose.position.x = arm.current_end_effector_frame.p.x();
    pose.position.y = arm.current_end_effector_frame.p.y();
    pose.position.z = arm.current_end_effector_frame.p.z();
    arm.current_end_effector_frame.M.GetQuaternion(
        pose.orientation.x, pose.orientation.y,
        pose.orientation.z, pose.orientation.w);
    pose_pub->publish(pose);

    geometry_msgs::msg::Twist twist;
    twist.linear.x = arm.current_end_effector_dot.vel.x();
    twist.linear.y = arm.current_end_effector_dot.vel.y();
    twist.linear.z = arm.current_end_effector_dot.vel.z();
    twist.angular.x = arm.current_end_effector_dot.rot.x();
    twist.angular.y = arm.current_end_effector_dot.rot.y();
    twist.angular.z = arm.current_end_effector_dot.rot.z();
    twist_pub->publish(twist);

    geometry_msgs::msg::Wrench wrench;
    wrench.force.x = arm.current_end_effector_wrench.force.x();
    wrench.force.y = arm.current_end_effector_wrench.force.y();
    wrench.force.z = arm.current_end_effector_wrench.force.z();
    wrench.torque.x = arm.current_end_effector_wrench.torque.x();
    wrench.torque.y = arm.current_end_effector_wrench.torque.y();
    wrench.torque.z = arm.current_end_effector_wrench.torque.z();
    wrench_pub->publish(wrench);
}
\end{lstlisting}

\textbf{验证方式：} Foxglove/RViz 能同时看到左右臂 joint state、末端 pose、末端 twist、估计 wrench。

\subsection{P2-2 把增益、限位、方向和串口参数移出硬编码}

\textbf{现象：} 当前电机限位、PID 增益、串口名 \verb|/dev/ttyACM0|、波特率等写死在 C++ 代码中。双臂调试时建议放到 YAML 参数文件。

\textbf{建议配置：}
\begin{lstlisting}
serial:
  port: /dev/ttyACM0
  baudrate: 921600

left_arm:
  motor_direction: [1, -1, 1, 1, 1, 1]
  pos_limit_min: [-1.5, -1.3, -1.5, -1.5, -1.5, -1.57]
  pos_limit_max: [1.5, 1.5, 1.5, 1.5, 1.5, 1.57]
  default_kp: [350, 300, 280, 10, 0, 0]
  default_kd: [1.9, 2.0, 2.0, 0.6, 0, 0]

right_arm:
  motor_direction: [1, -1, 1, 1, 1, 1]
  pos_limit_min: [-1.5, -1.3, -1.5, -1.5, -1.5, -1.57]
  pos_limit_max: [1.5, 1.5, 1.5, 1.5, 1.5, 1.57]
\end{lstlisting}

\textbf{建议代码：}
\begin{lstlisting}[language=C++]
this->declare_parameter<std::string>("serial.port", "/dev/ttyACM0");
this->declare_parameter<int>("serial.baudrate", 921600);

const auto port = this->get_parameter("serial.port").as_string();
const auto baud = this->get_parameter("serial.baudrate").as_int();
\end{lstlisting}

\subsection{P2-3 增加通信超时 watchdog}

\textbf{现象：} 当前下位机只根据 \verb|Is_USBget| 区分是否收到数据，但缺少持续超时失能策略。建议上位机和下位机都补 watchdog。

\textbf{下位机建议代码：}
\begin{lstlisting}[language=C]
static uint32_t last_usb_tick = 0;

// CDC_Receive_HS 中:
last_usb_tick = HAL_GetTick();
Is_USBget = 1;

// main while 中:
if ((HAL_GetTick() - last_usb_tick) > 100) {
    fdcanx_send_data(&hfdcan1, 0x01, disable_data, 8);
    fdcanx_send_data(&hfdcan1, 0x02, disable_data, 8);
    fdcanx_send_data(&hfdcan1, 0x03, disable_data, 8);
    fdcanx_send_data(&hfdcan1, 0x04, disable_data, 8);
    fdcanx_send_data(&hfdcan1, 0x05, disable_data, 8);
    fdcanx_send_data(&hfdcan1, 0x06, disable_data, 8);
    fdcanx_send_data(&hfdcan2, 0x01, disable_data, 8);
    fdcanx_send_data(&hfdcan2, 0x02, disable_data, 8);
    fdcanx_send_data(&hfdcan2, 0x03, disable_data, 8);
    fdcanx_send_data(&hfdcan2, 0x04, disable_data, 8);
    fdcanx_send_data(&hfdcan2, 0x05, disable_data, 8);
    fdcanx_send_data(&hfdcan2, 0x06, disable_data, 8);
}
\end{lstlisting}

\subsection{P2-4 修正拼写和命令兼容}

\textbf{现象：} 代码中存在 \verb|Impdence|、\verb|laod|、\verb|clearlaod| 等拼写。建议保留旧拼写兼容，但新接口使用正确命名。

\textbf{建议代码：}
\begin{lstlisting}[language=C++]
if (cmd == "Impedance" || cmd == "Impdence") {
    return State::Impdence; // 内部枚举可后续统一重命名
}

if (cmd == "load" || cmd == "laod") {
    llarm6dof.cmd = "load";
}

if (cmd == "clearload" || cmd == "clearlaod") {
    llarm6dof.cmd = "clearload";
}
\end{lstlisting}

\section{P3：工程化完善}

\subsection{P3-1 增加协议单元测试}

\textbf{目的：} 防止 112 字节协议、CRC、左右臂槽位在后续改动中被破坏。

\textbf{建议测试伪代码：}
\begin{lstlisting}[language=C++]
TEST(ProtocolPacking, DualArmPayloadLayout)
{
    MotorControl motor(serial);
    MotorControl::ArmCommand left;
    MotorControl::ArmCommand right;

    left.pos[0] = 0.1f;
    right.pos[0] = -0.2f;
    left.gripper = 10;
    right.gripper = 20;

    auto payload = motor.packDualArmPayloadForTest(left, right);

    ASSERT_EQ(payload.size(), 112);
    EXPECT_EQ(payload[48], 10);
    EXPECT_EQ(payload[104], 20);
    EXPECT_NE(payload[0], payload[56]);
}
\end{lstlisting}

\subsection{P3-2 建立仿真/无硬件模式}

\textbf{目的：} 在没有 H7 和电机时也能测试 ROS2 话题、状态机和规划逻辑。

\textbf{建议代码：}
\begin{lstlisting}[language=C++]
class ITransport {
public:
    virtual ~ITransport() = default;
    virtual void send(const std::vector<uint8_t> &frame) = 0;
    virtual std::vector<uint8_t> receive() = 0;
};

class SerialTransport final : public ITransport { /* real serial */ };
class SimTransport final : public ITransport {
public:
    void send(const std::vector<uint8_t> &frame) override { last_frame_ = frame; }
    std::vector<uint8_t> receive() override { return simulated_feedback_; }
private:
    std::vector<uint8_t> last_frame_;
    std::vector<uint8_t> simulated_feedback_;
};
\end{lstlisting}

\subsection{P3-3 补齐文档和启动文件}

\textbf{建议补充：}
\begin{itemize}[leftmargin=2em]
  \item \verb|bringup_dual_arm.launch.py|：启动主控制节点、键鼠节点、foxglove bridge。
  \item \verb|config/dual_arm.yaml|：集中配置串口、左右臂参数、限位、默认模式。
  \item \verb|docs/protocol.md|：说明 112 字节 payload 的槽位布局。
  \item \verb|docs/safety.md|：说明急停、零位校准、上电顺序、调试禁忌。
\end{itemize}

\section{推荐执行顺序}

\begin{enumerate}[leftmargin=2em]
  \item 先完成 P0-1 和 P0-2，修复下位机缓存和右臂下发。这是双臂实机安全底线。
  \item 再完成 P0-3 和 P0-4，确保保护和状态切换不会留下危险状态。
  \item 接着做 P1-1 和 P1-2，把上位机从 6 轴控制真正升级为 12 轴控制。
  \item 然后做 P1-3、P1-4、P1-5，统一双臂话题、命令和轨迹来源。
  \item 主链路稳定后做 P2/P3，提高调试效率、可观测性和长期维护质量。
\end{enumerate}

\section{最小闭环验收标准}

完成 P0 和 P1 后，建议按以下标准验收：

\begin{itemize}[leftmargin=2em]
  \item 上位机能分别发布左臂/右臂关节目标，另一侧不受影响。
  \item 下位机两条 FDCAN 总线收到不同的 6 电机指令。
  \item 任意一侧位置、速度、力矩越界都会触发失能。
  \item 串口断开或上位机停止发送超过设定时间，下位机进入安全态。
  \item Foxglove/RViz 能看到左右臂 joint states 和末端位姿。
\end{itemize}

\end{document}
